% samosa: A: n×n; B: n×n; u: n; v: n; x: n
\documentclass{article}
\title{Three identities of the transpose and the inverse}
\author{A. Author}
\begin{document}
\begin{abstract}
We record three identities that hold at every size, and we check each step
of the derivations exactly.
\end{abstract}

\section{Introduction}
Let $A$ and $B$ be square matrices of the same size, and let $u$ and $v$ be
column vectors. We write $x$ for a column vector as well.

\section{The transpose of a product}
\begin{theorem}
The transpose of a product keeps the order.
\end{theorem}
\begin{proof}
\begin{align*}
(A B)^\top &= B^\top A^\top
\end{align*}
\end{proof}

\section{A transpose twice}
\begin{theorem}
A transpose twice is nothing.
\end{theorem}
\begin{proof}
\begin{align*}
(A^\top)^\top &= A
\end{align*}
\end{proof}

\section{A rank one update}
\begin{theorem}
The inverse of a rank one update.
\end{theorem}
\begin{proof}
\begin{align*}
(A + u v^\top)^{-1} &= A^{-1} - \frac{A^{-1} u v^\top A^{-1}}{1 + v^\top A^{-1} u}
\end{align*}
\end{proof}

\end{document}
